% Plantilla de pandoc para generar informe.Rmd con la plantilla de informes UC
% (plantilla/style.cls). La usa el campo output: pdf_document: template de informe.Rmd.
\documentclass{plantilla/style}

% --- Definiciones que pandoc necesita ---------------------------------------
\makeatletter
\newsavebox\pandoc@box
\newcommand*\pandocbounded[1]{%
  \sbox\pandoc@box{#1}%
  \Gscale@div\@tempa{\textheight}{\dimexpr\ht\pandoc@box+\dp\pandoc@box\relax}%
  \Gscale@div\@tempb{\linewidth}{\wd\pandoc@box}%
  \ifdim\@tempb\p@<\@tempa\p@\let\@tempa\@tempb\fi
  \ifdim\@tempa\p@<\p@\scalebox{\@tempa}{\usebox\pandoc@box}%
  \else\usebox{\pandoc@box}%
  \fi}
\makeatother
\providecommand{\passthrough}[1]{#1}
$if(numbersections)$
\setcounter{secnumdepth}{3}
$else$
\setcounter{secnumdepth}{-\maxdimen}
$endif$
\graphicspath{{./}}
\renewcommand{\contentsname}{Índice de contenido}
\renewcommand{\listfigurename}{Índice de figuras}
\renewcommand{\listtablename}{Índice de tablas}
\addto\captionsspanish{\renewcommand{\tablename}{Tabla}\renewcommand{\figurename}{Figura}}
$for(header-includes)$
$header-includes$
$endfor$

\begin{document}

% --- Encabezado de cada página ----------------------------------------------
\fancyhead[L]{%
  \begin{picture}(0,0) \put(-40,26){\includegraphics[width=15mm]{plantilla/img/logo-uc-3.pdf}} \end{picture}
  \put(8,50){\textcolor{gray}{\scriptsize{\begin{tabular}{l}
    PONTIFICIA UNIVERSIDAD CATÓLICA DE CHILE\\
    FACULTAD DE MATEMÁTICAS\\
    DEPARTAMENTO DE ESTADÍSTICA
  \end{tabular}}}}
}
\rhead{}

% --- Portada -----------------------------------------------------------------
\begin{titlepage}
\includegraphics[width=2cm]{plantilla/img/logo-uc-4.pdf}
\vspace*{-2.2cm}

\begin{tabular}{l}
\hspace{2cm} PONTIFICIA UNIVERSIDAD CATÓLICA DE CHILE\\
\hspace{2cm} FACULTAD DE MATEMÁTICAS\\
\hspace{2cm} DEPARTAMENTO DE ESTADÍSTICA\\
\end{tabular}

\begin{center}
\vspace{3cm}
{\scshape\LARGE \textbf{$title$} \par}
\rule{80mm}{0.1mm}

{\itshape\Large $subtitle$ \par}
\vfill
{\Large \textbf{$grupo$} \par}
\vspace{0.5cm}
{\large \textbf{Integrantes:}\\\normalsize $for(author)$$author$$sep$, $endfor$ \par}

\medskip
\textbf{Fecha de entrega:} $date$
\vspace{3cm}
\end{center}
\end{titlepage}
\newpage

% Sin índices: el informe es breve y la numeración de secciones basta (adaptación).
\setstretch{1}

$body$

\end{document}
